\section{Algorithm}
\label{sec:algorithm}

\subsection{Notation and Setup}
\label{sec:alg:notation}

Let $G = (V, E, w)$ be the census-tract adjacency graph for a state (or
sub-region), where $V$ is the set of census tracts, $E$ the set of
shared-boundary edges, and $w: E \to \mathbb{R}_{>0}$ the TIGER boundary length.
Let $\text{pop}: V \to \mathbb{Z}_{>0}$ be census population.
Let $k \geq 2$ be the number of districts to produce.
The bisection tree for seat count $k$ is determined by the split prescription
of \citet{dellaLibera2026apportion}: at each node with $k'$ target districts,
the node applies a $p'$-way split where $p' = \text{largest prime factor of } k'$
(or $p' = 2$ for prime $k' > 3$, using floor/ceil targets).

BisectionEnsemble replaces the METIS call at each bisection node (where $p' = 2$)
with a local ReCom ensemble. It is defined by two hyperparameters:
\begin{itemize}
  \item $p \in [0, 1]$: the percentile of the edge-cut distribution to select.
  \item $T \in \mathbb{Z}_{>0}$: the number of ReCom steps to run.
\end{itemize}

\subsection{The BisectionEnsemble Primitive}
\label{sec:alg:primitive}

\begin{definition}[Local bisection ensemble]
  \label{def:be-primitive}
  Given a subgraph $H = (V_H, E_H, w)$ with populations $\text{pop}|_{V_H}$ and
  target district count $k' = 2$ (with population targets $t_1 = \lfloor k/2
  \rfloor$ and $t_2 = \lceil k/2 \rceil$ as fractions of the node's total
  population), \textup{BisectionEnsemble}$(H, p, T)$ proceeds as follows:
  \begin{enumerate}
    \item \textbf{Initialise}: Run METIS 2-way partitioning on $H$ with targets
          $(t_1, t_2)$ to obtain an initial bisection $\pi_0 = (A_0, B_0)$.
    \item \textbf{ReCom chain}: For $t = 1, \ldots, T$:
          \begin{enumerate}
            \item Find a random boundary edge $(u, v) \in E_H$ with
                  $u \in A_{t-1}$, $v \in B_{t-1}$.
            \item Form the merged subgraph $M = H[A_{t-1} \cup B_{t-1}] = H$.
            \item Sample a spanning tree $\mathcal{T} \sim \text{UST}(M)$ via
                  Wilson's algorithm.
            \item Find a cut edge $e^* \in \mathcal{T}$ such that removing $e^*$
                  produces two components $(A', B')$ with
                  $|\text{pop}(A') - t_1 \cdot \text{pop}(H)| \leq \varepsilon$.
            \item If no such $e^*$ exists, reject (keep $\pi_{t-1}$); else
                  accept: $\pi_t = (A', B')$.
          \end{enumerate}
    \item \textbf{Collect}: Let $\Pi = \{\pi_t : \pi_t \text{ was accepted}\}$
          with $|\Pi| = m$. For each $\pi \in \Pi$, compute the edge-cut
          weight $\text{EC}(\pi) = \sum_{(u,v) \in \delta(\pi)} w(u,v)$.
    \item \textbf{Select}: Sort $\Pi$ by $\text{EC}(\cdot)$. Return
          $\pi^* = \Pi[\lfloor p \cdot m \rfloor]$ (zero-indexed).
  \end{enumerate}
\end{definition}

\noindent\textbf{Remark on step 2(c).} In BisectionEnsemble, the full merge
region $M$ is simply the subgraph $H$ itself (since we are splitting $H$ into
exactly 2 parts). The spanning tree is thus sampled over the complete local
subgraph at each step. This differs from a general $k$-way ReCom step where
only two of the $k$ districts are merged; here, the entire node region is
always the merge domain.

\subsection{Integration with the Bisection Tree}
\label{sec:alg:tree}

Algorithm~\ref{alg:bisection-ensemble} gives the full recursive procedure.
BisectionEnsemble replaces every 2-way METIS call in the bisection tree.
For $p'$-way splits with $p' > 2$ (ApportionRegions nodes), standard METIS
is used (BisectionEnsemble is not applicable to $p' > 2$ nodes since those
are not bisections). A pure binary tree (all seat counts $k = 2^d$) applies
BisectionEnsemble at every node.

\begin{algorithm}[h]
\caption{BisectionEnsemble redistricting}
\label{alg:bisection-ensemble}
\begin{algorithmic}[1]
\Require Subgraph $H$, seat count $k$, percentile $p$, steps $T$
\Ensure Set of $k$ district subgraphs
\If{$k = 1$}
  \Return $\{H\}$
\EndIf
\State $(p', k_1, \ldots, k_{p'}) \leftarrow \text{split}(k)$ \Comment{Per AR split prescription}
\If{$p' = 2$} \Comment{Bisection node: use local ensemble}
  \State $(A, B) \leftarrow \textup{BisectionEnsemble}(H, p, T)$
\Else \Comment{$p' > 2$ node: use METIS directly}
  \State $(R_1, \ldots, R_{p'}) \leftarrow \text{METIS}(H, p')$
\EndIf
\For{each sub-region $R_i$ with target $k_i$}
  \State $D_i \leftarrow \textup{BisectionEnsemble-Redist}(R_i, k_i, p, T)$
\EndFor
\Return $\bigcup_i D_i$
\end{algorithmic}
\end{algorithm}

\subsection{Key Theoretical Properties}
\label{sec:alg:theory}

\begin{theorem}[Tractability]
  \label{thm:tractable}
  BisectionEnsemble never encounters a bipartition failure. That is, for any
  local subgraph $H$ and any $T \geq 1$, the ReCom chain in Definition~\ref{def:be-primitive}
  is well-defined and will produce at least one accepted plan within
  $O(|V_H|)$ steps in expectation.
\end{theorem}

\begin{proof}
  At each step $t$, the merge domain is $M = H$, which has a uniform spanning
  tree over $V_H$. The population balance constraint requires finding a cut
  edge $e^*$ such that the two components have populations within $\varepsilon$
  of the targets $(t_1, t_2) = (1/2, 1/2)$ (for equal split). By the
  Aldous-Broder theorem \citep{aldous1990random}, the uniform spanning tree
  of $H$ is a connected tree on $V_H$. Removing any edge $e \in \mathcal{T}$
  produces two connected components; as $e$ ranges over all $|V_H| - 1$ edges,
  the components $(A_e, B_e)$ form a family of nested cuts. For any target
  ratio $r = t_1 / (t_1 + t_2) \in (0, 1)$ and tolerance $\varepsilon > 0$,
  the intermediate-value property of the discrete population function guarantees
  the existence of at least one edge $e^*$ with population balance within
  $\varepsilon$ of $r$, provided $\varepsilon \geq \max_{v \in V_H} \text{pop}(v)
  / \text{pop}(H)$ (i.e., the tolerance is at least one tract's population share).
  Since BisectionEnsemble uses population balance tolerance $\varepsilon \geq 0.5\%$
  (statutory tolerance) and tracts represent $\ll 1\%$ of any node's population,
  the acceptance condition is always satisfiable. The chain therefore accepts
  at least one plan in finite steps.
\end{proof}

\begin{proposition}[Contiguity preservation]
  \label{prop:contiguity}
  Every plan returned by BisectionEnsemble is contiguous. Specifically, every
  district subgraph produced at every tree node is a connected induced subgraph
  of the census-tract adjacency graph $G$.
\end{proposition}

\begin{proof}
  The spanning tree $\mathcal{T}$ of the merge domain $M = H$ is a connected tree.
  Removing any single edge $e^*$ from $\mathcal{T}$ produces exactly two
  connected subtrees $(A, B)$, both of which are connected. Since $A$ and $B$
  are subgraphs of $H \subseteq G$, they are connected induced subgraphs of $G$.
  Inductively, all sub-regions passed to recursive calls are connected, and the
  final districts are connected.
\end{proof}

\begin{proposition}[Complexity]
  \label{prop:complexity}
  For a bisection tree of depth $d = O(\log k)$ and a state graph with $n$ tracts,
  BisectionEnsemble runs in $O(T \cdot n \log n)$ total time across all tree nodes.
\end{proposition}

\begin{proof}
  At depth $\ell$ of a binary tree, there are $2^\ell$ nodes, each managing a
  subgraph of $\approx n / 2^\ell$ tracts. Each node runs $T$ ReCom steps.
  Wilson's algorithm samples a uniform spanning tree in $O(|V_H|)$ expected
  time \citep{wilson1996generating}. The cost per node is $O(T \cdot n / 2^\ell)$.
  Summing over all $d$ levels:
  \[
    \sum_{\ell=0}^{d-1} 2^\ell \cdot O\!\left(\frac{Tn}{2^\ell}\right)
    = \sum_{\ell=0}^{d-1} O(Tn)
    = O(Tn \cdot d)
    = O(Tn \log k).
  \]
  Since $n \log k \leq n \log n$, the total time is $O(T n \log n)$.
\end{proof}

\subsection{Parallelism}
\label{sec:alg:parallel}

Tree nodes at the same depth are independent: each node's ensemble depends
only on the subgraph it inherits from its parent, which is fixed once the
parent's bisection is chosen. This gives two natural parallelism strategies:
\begin{itemize}
  \item \textbf{Level-parallel}: Execute all nodes at depth $\ell$ in parallel,
        then proceed to depth $\ell + 1$. Requires synchronisation at level
        boundaries but achieves $2^\ell$-fold parallelism at level $\ell$.
  \item \textbf{Asynchronous}: Begin all nodes immediately; sub-regions start
        as soon as their parent's bisection is selected. This is a pipeline
        over the tree depth with $d$ concurrent stages.
\end{itemize}
The \texttt{redist-cli} implementation uses Rayon's work-stealing thread pool,
which naturally implements the asynchronous strategy: nodes at deeper levels
are spawned as Rayon tasks as soon as their parent completes.
